\newif \ifshare
% \sharetrue % Comment this if you want animation
\ifshare % "Share mode" without animation.
\documentclass[table, trans, aspectratio = 169]{beamer}
\else % "Presentation mode" with animation.
\documentclass[table, aspectratio = 169]{beamer}
\fi
\usepackage{default}
\usepackage[T1]{fontenc}
\usepackage{color}
\usepackage{graphicx}
\usepackage{subfig}
\usepackage{diagbox}


\let\pgfmathMod=\pgfmathmod\relax

\graphicspath{{../figs/PDF/}}

\usetheme{Boadilla}

\definecolor{struct}{HTML}{03a9f4}
\definecolor{alert}{HTML}{f44336}
\definecolor{example}{HTML}{aeea00}
\definecolor{good}{HTML}{8bc34a}
\definecolor{notgoodnotbad}{HTML}{ff9800}

\setbeamercolor{structure}{fg = struct}
\setbeamercolor{normal text}{fg = white, bg = black}
\setbeamercolor{example text}{fg = example}
\setbeamercolor{alerted text}{fg = alert}
\setbeamercolor{footline}{fg = white}

\title[eBPF ready kernel 5.10 for \texttt{minikube} (Rejekt 2022)]{eBPF Ready Kernel 5.10 for \texttt{minikube}}
\subtitle{Rejekt 2022}
\author[Francis Laniel (\texttt{flaniel@linux.microsoft.com})]{Francis Laniel\\\texttt{flaniel@linux.microsoft.com}}
\date{15 May 2022}

% Custom title page.
\defbeamertemplate*{title page}{customized}[1][]{
	\centering
	\usebeamerfont{title}\usebeamercolor[fg]{title}\inserttitle\par
	\usebeamerfont{subtitle}\usebeamercolor[fg]{subtitle}\insertsubtitle\par
	\bigskip
	\usebeamerfont{author}\usebeamercolor[fg]{normal text}\textbf{\insertauthor}\par
	\bigskip
	\usebeamerfont{date}\usebeamercolor[fg]{normal text}\textbf{\insertdate}\par
	\bigskip
	\bigskip

	\begin{columns}
		\begin{column}{.5\textwidth}
			\centering

			\includegraphics[scale=2]{microsoft.pdf}
		\end{column}
		\begin{column}{.5\textwidth}
			\centering

			\includegraphics{kinvolk.pdf}
		\end{column}
	\end{columns}
}

\begin{document}
	% Put these parameters here to avoid compilation error:
	% "! LaTeX Error: Missing \begin{document}."
	% Remove the navigation bar, this is useless...
	\setbeamertemplate{navigation symbols}{}
	% Use square instead of bubbles, see:
	% https://tex.stackexchange.com/a/69721
	\setbeamertemplate{section in toc}[square]
	% Modify the shaded value to 60% instead of 20%, see:
	% https://tex.stackexchange.com/a/66703
	\setbeamertemplate{sections/subsections in toc shaded}[default][50]
	% Use circle instead of bubbles for itemize, see:
	% \setbeamertemplate{itemize items}[circle]
	\setbeamertemplate{itemize items}[square]
	\setbeamertemplate{enumerate items}[square]

	\maketitle

	\section{\texttt{minikube}}
	\begin{frame}
		\frametitle{\texttt{minikube}}
		\framesubtitle{What it is?}

		\texttt{minikube} implements \cite{minikube_contributors_minikube_nodate}:
		\begin{quote}
			[...] a local Kubernetes cluster [...]. minikube's primary goals are to be the best tool for local Kubernetes application development and to support all Kubernetes features that fit.
		\end{quote}

		\bigskip

		\onslide<2>{
			To deploy local Kubernetes, \texttt{minikube} offers three drivers \cite{minikube_contributors_drivers_2021}:
			\begin{itemize}
				\item \texttt{none}
				\item \texttt{docker}
				\item \textcolor{struct}{\texttt{kvm}}
			\end{itemize}
		}
	\end{frame}

	\section{Problem and context}
	\begin{frame}[fragile]
		\frametitle{Problem and context \cite{starovoitov_patch_2020}}

		\begin{columns}
			\begin{column}{.5\textwidth}<1->
				\centering

				\includegraphics[scale = .33]{kernel_5.7.pdf}
			\end{column}
			\begin{column}{.5\textwidth}<2->
				\centering

				\includegraphics[scale = .33]{kernel_minikube.pdf}
			\end{column}
		\end{columns}

		\bigskip

		\onslide<3->{
			\texttt{minikube} ISO is built using \texttt{buildroot} which is \cite{buildroot_contributors_buildroot_2022}:
			\begin{quote}
				[...] a simple, efficient and easy-to-use tool to generate embedded Linux systems through cross-compilation.
			\end{quote}

			It can also be used to generate "virtualization ready" image with, \textit{e.g.}, \texttt{qemu\_x86\_64\_defconfig} \cite{zacarias_buildroot_2011}.
		}
	\end{frame}

	\section{Contribution}
	\begin{frame}
		\frametitle{Contribution}
		\framesubtitle{Bumping kernel version}

		\centering

		% CONFIG_TMPFS to be able to use mount -t tmpfs in init.
		% CONFIG_PCI to have network.
		% CONFIG_BRIDGE_NETFILTER to have
		% /proc/sys/net/bridge/bridge-nf-call-iptables
		% which is needed by kubeadm.

		\includegraphics<1>[scale=2]{buildroot_process-fig1.pdf}%
		\includegraphics<2>[scale=2]{buildroot_process-fig2.pdf}%
		\includegraphics<3>[scale=2]{buildroot_process-fig3.pdf}%
		\includegraphics<4>[scale=2]{buildroot_process-fig4.pdf}%
		\includegraphics<5>[scale=2]{buildroot_process-fig5.pdf}%
		\includegraphics<6>[scale=2]{buildroot_process-fig6.pdf}%
	\end{frame}

	\begin{frame}
		\frametitle{Contribution}
		\framesubtitle{Adding eBPF needed \texttt{CONFIG\_*}}

		\centering

		\includegraphics<1>[scale=2]{buildroot_epbf_process-fig1.pdf}%
		\includegraphics<2>[scale=2]{buildroot_epbf_process-fig2.pdf}%
	\end{frame}

	\section{Conclusion}
	\begin{frame}
		\frametitle{Conclusion and future work}

		Conclusion:
		\begin{enumerate}
			\item eBPF ready kernel 5.10 minikube image is available \cite{minikube_contributors_minikube_2021}!
			\item \texttt{sysdig} bug was fixed in upstream buildroot \cite{laniel_buildroot_2022}!
			\item \texttt{pahole} package was added to upstream buildroot \cite{laniel_buildroot_2022-1}!
		\end{enumerate}

		\bigskip

		\onslide<2>{
			Future work:
			\begin{enumerate}
				\item Merge this work in upstream minikube \cite{laniel_use_2021}.
				\item Keep up to date upstream buildroot packages.
			\end{enumerate}
		}
	\end{frame}

	\begin{frame}[allowframebreaks, noframenumbering]
		\frametitle{Bibliography}
		\setbeamertemplate{bibliography item}[text]

		\begin{scriptsize}
			\bibliographystyle{IEEEtran}
			\bibliography{beamer}

			\bigskip
			Some pictures from this presentation contain pictures from the following authors: \href{https://github.com/Templarian/MaterialDesign}{Austin Andrew}, \href{https://github.com/bytedance/IconPark}{ByteDance}, \href{https://github.com/EmojiTwo/emojitwo}{Emoji One}, \href{https://github.com/ant-design/ant-design-icons}{HeskeyBaozi} and \href{https://github.com/icons8/line-awesome}{Icons8}
		\end{scriptsize}
	\end{frame}
\end{document}
